\section{Related Work and Limits of the Claim}
\label{sec:related}

\paragraph{Workspaces and agent-native file systems.}
OverlayFS supplies the layered semantics that motivate the private upper
and shared lower~\cite{overlayfs}. BranchFS provides copy-on-write branch
contexts for agentic exploration, including a commit/abort lifecycle; its
workshop prototype and proposed kernel abstraction must be
distinguished~\cite{branchfs2026}. YoloFS studies staging, snapshots, and
permissions to improve visibility and control over agent
effects~\cite{yolofs2026}. AgentFS also documents a file-system and tool-state
abstraction for agents~\cite{agentfs2025}. These efforts make an
``agent filesystem'' label insufficient as a novelty claim. \sys instead
targets the costs of shared remote lower metadata and object access. It does
not replace process isolation or provide all of these systems' user-control
and commit semantics.

Table~\ref{tab:comparison} makes the intended advantage more precise. The
last column is a hypothesis to be tested under matched persistence and cache
budgets, rather than a claim that the prototype dominates every row.

\begin{table*}[t]
\centering
\scriptsize
\begin{tabularx}{\textwidth}{@{}p{1.60cm}p{1.95cm}p{2.15cm}p{2.0cm}Y@{}}
\toprule
System & Shared baseline & Private view & Remote metadata/data & Boundary tested in this work\\
\midrule
OverlayFS & local lower tree & kernel upper/copy-up & none by itself & remote immutable lower, explicit object reads, and durable upper\\
DeltaFS & immutable logs and manifests & per-job namespace & object store plus registry & current-value upper, directory-oriented pages, and online POSIX lookup\\
BranchFS & local COW branch context & branch commit/abort & prototype-local & remote persistence and multi-client lower sharing\\
CVMFS/Nydus & lazy immutable catalog/image & requires an external upper & cached objects/blobs & writable workspace semantics and metadata/data locator co-location\\
JuiceFS & chunk/slice/block volume & shared mutable POSIX view & object store plus KV & lower-query independence and an explicit persisted $(B,q)$ layout\\
\sys & selected packed lower & tombstone-aware private upper & packed objects plus control KV & sparse mutation, cold lower access, and layout effects in one measured contract\\
\bottomrule
\end{tabularx}
\caption{Comparison axes used to choose baselines. Existing systems already
provide individual mechanisms; the claimed opportunity is their scoped
composition for remote agent workspaces.}
\label{tab:comparison}
\end{table*}

\paragraph{Namespace and metadata organization.}
TableFS, IndexFS, LocoFS, and InfiniFS demonstrate that record organization,
partitioning, and caching can substantially change metadata access
paths~\cite{tablefs2013,indexfs2014,locofs2017,infinifs2022}.
Recent systems also revisit hierarchical object metadata and distributed
update coordination~\cite{mantle2025,switchfs2026}. We compare mechanisms at
their actual service boundaries instead of assuming that every baseline
performs a remote per-file lookup.

DeltaFS is a particularly close conceptual predecessor. It uses private
namespaces and immutable mutation logs over shared storage, with selective
composition of prior state~\cite{deltafs2021}. \sys cannot claim private
deltas or object-resident snapshots as inventions. Its proposed distinction
is a directory-oriented lower layout that combines common attributes and
data locators in paged access units, composed with an online private upper
and a backend layout interface. Whether this particular combination is
worth its complexity is an evaluation question.

\paragraph{Immutable distribution and storage backends.}
CernVM-FS provides demand-driven software distribution with indexed
catalogs, while Nydus describes metadata/data organization for lazy
image access~\cite{cvmfs2011,nydus_rafs}. Slacker and DADI also address
image distribution and startup costs~\cite{slacker2016,dadi2020}.
These are stronger points of comparison for immutable baselines than a
model that downloads every file before use. The experiment plan therefore
requires a mature lazy-lower baseline and counts image construction and
staging separately.

Distributed systems such as Ceph and Tectonic separate aspects of metadata
and data services~\cite{ceph2006,tectonic2021}; this separation itself is not
a contribution of \sysname. Erasure-coded storage research considers
repair cost, changing redundancy over a file's lifetime, and the relation
between data placement and redundancy~\cite{azureec2012,morph2024,okapi2025}.
Okapi's separation of striping and redundancy grouping is different from
our $(B,q)$ distinction, although both motivate making coupled decisions
explicit. We keep code rate fixed and measure checksum, tail, and repair
cost rather than inferring a lower redundancy ratio from smaller blocks.

\paragraph{Scope and unresolved evidence.}
The design is most plausible for stable shared baselines and sparse private
changes. A workload with poor name locality, high mutation, or rapid
snapshot churn may favor existing mutable metadata services or local
overlays. The current Entity mount limits scale and lifecycle guarantees;
the current automatic B policy can be identical to the default, and its
upper objects do not exercise the above-1-MiB branch through ordinary
FUSE writes. These are material limits, not editorial details. Results
must determine whether all three title mechanisms deserve equally strong
performance claims.
